\section{Related Work}

\paragraph{Multilingual Speech Datasets.} 
Current multilingual speech datasets with transcriptions are drawn from a variety of domains, including Wikipedia~\citep{ardila2019common,conneau2022fleurs}, political speech~\citep{wang2021voxpopuli}, audiobooks~\citep{pratap2020mls} to name a few.  
They are limited to about 100 languages with some datasets only containing data for European languages \citep{wang2021voxpopuli,pratap2020mls}.
Multilingual datasets without transcriptions include VoxLingua107~\citep{valk2020voxlingua} spanning data in 107 languages, or VoxPopuli~\citep{wang2021voxpopuli} which contains large amounts of unlabeled data for European languages.
\citet{leong2022bloom} covers 56 low-resource languages and 428 hours of data.

Compared to prior work utilizing read versions of the New Testament~\cite{black2019cmu}, MMS covers many more languages (58\% more for \MMS{} and over five times more for \GRN{}) and our alignments are of much higher quality which leads to better models as we show in~\textsection\ref{sec:data-comp-cmu}.
We also use the resulting data to train self-supervised models, build speech recognition systems as well as language identification models while as~\citet{black2019cmu} focused on speech synthesis.



\paragraph{Multilingual Automatic Speech Recognition (ASR).}
Prior work on multilingual speech recognition includes both non-neural methods~\citep{burget2010multilingual,lin2009study}, hybrid neural and HMM models~\citep{heigold2013multilingual} and more recently neural systems~\citep{cho2018multilingual,toshniwal2018multilingual,kannan2019mult,li2019bytes}.
\citet{boli2021scaling} built multilingual ASR for up to 15 languages, \citet{Pratap2020multiasr,lugosch2022icassp, tjandra2022massively} trained and explored several strategies for 50\texttt{+} languages. 

Whisper~\citep{radford2022whisper} mines 680K hours of data from the web and their model supports the transcription of 99 different languages. 
\cite{zhang2023usm} trained a multilingual ASR model based on YouTube audio data and successfully scaled to 100 languages. 
A notable exception is~\citet{li2022asr2k} who created ASR for 1,909 languages by mapping the phoneme-like output of an eight language multilingual model to appropriate phonemes for the language of interest. 
In contrast, MMS uses actual paired speech and text for over 1,100 languages and present a comparison to their approach below (\textsection\ref{sec:data-comp-asr2k}).



\paragraph{Spoken Language Identification (LID).}
Whisper~\citep{radford2022whisper} also supports language identification and can distinguish between 99 different languages. 
\citet{fan2021exploring} utilizes wav2vec 2.0 for language identification using the API17-OLR dataset that consists of ten Asian languages. 
Later, \citet{tjandra2022improved} demonstrated that a cross-lingual self-supervised model can improve language identification performance by training a language identification model for 26 languages using a proprietary dataset. 
\citet{babu22_interspeech} fine-tunes a pre-trained model to perform LID for 107 languages using the VoxLingua-107 dataset~\citep{valk2020voxlingua}.  
In this work, we scale the number of languages to over 4,000 which to our knowledge is the broadest coverage spoken language identification model so far.


\paragraph{Multilingual Text-To-Speech (TTS).} 
Speech synthesis has been undergoing a transition from controlled settings to the generation of more diverse speech, with multilinguality being a crucial aspect~\citep{casanova2022yourtts,zhang2023speak}. 
However, the lack of multilingual training data, particularly for low-resource languages, presents a common obstacle in scaling TTS to more languages.
To address data scarcity, prior work explored various approaches including byte encoding to unify text representations which was evaluated on English, Spanish, and Chinese~\citep{li2019bytes}
Further studies explored other input representations, including phonemes~\citep{zhang2019learning} and phonological features \citep{Staib2020Phonological}. 

Additionally, various modeling schemes have been developed to encourage knowledge sharing between languages, such as parameter generation networks~\citep{Nekvinda2020one} and leveraging unpaired speech or text data for pre-training~\citep{Saeki2023Virtuoso,Saeki2023learning}.
Despite these efforts, most prior work still covers a small number of languages but there are a few efforts which scaled to 46 languages~\citep{he2021multi} or use unsupervised techniques to scale to 101 languages~\citep{Saeki2023Virtuoso}. 
\citet{meyer2022bibletts} also builds VITS models based on readings of the Bible but their work is limited to ten African languages.


\paragraph{Multilingual NLP.}
Multilinguality has been a very active research area in NLP where researchers introduced cross-lingually pre-trained sentence encoders~\citep{conneau2019cross} spanning 100 languages, or pre-trained multilingual sequence to sequence models~\citep{liu2020multilingual} applied to machine translation. 
Multilingual machine translation models have also been scaled to 200 languages~\citep{nllbteam2022language} and even 1,000 languages~\citep{bapna2022mt1k}.
